\documentclass{cuitthesis}

% ===== 论文信息：按实际内容填写 =====
\cuitclassification{}
\cuitudc{}
\cuitsecurity{公开}
\cuitstudentnumber{}
\cuitcoverctitle{论 文 中 文 题 目}
\cuitcoveretitle{Dissertation in English}
\cuitctitle{学位论文题目中文名称}
\cuitetitle{Dissertation in English}
\cuitauthor{}
\cuitcollege{}
\cuitdegreetype{}
\cuitstudytype{}
\cuitdiscipline{}
\cuitresearchfield{}
\cuitsupervisor{姓名\quad 职称}
\cuitcosupervisor{姓名\quad 职称}
\cuitfinishdate{\qquad 年\qquad 月}

\begin{document}

\makecuitcover

\cuitfrontmatter
\makecuitstatements

\begin{cuitabstractcn}
本文是在中国气象局公益性气象行业专项（项目编号：GYHY201006045）的资助下完成的。论文通过对人工与超声波雪深传感器、人工与激光雪深传感器以及超声波雪深传感器与激光雪深传感器的试验数据对比分析，对超声波雪深传感器和激光雪深传感器的测量方法做了系统的理论分析和深入的性能分析研究。本行业专项为了研究超声波雪深传感器和激光雪深传感器的测量特性，设计了室内实验和外场对比观测试验，对比超声波雪深传感器和激光雪深传感器的测量性能，同时对雪深传感器测量的数据做了系统的分析与研究，分析其数据与天气过程（风速、温度、降雪和融雪过程等）之间的关系以及雪深与雪量的关系。概括起来本文研究的主要内容如下所示：

(1) 设计实验室检测方案，通过在实验室内对传感器的静态检测，得到超声波雪深传感器和激光雪深传感器测量的数据，分析传感器的一致性和稳定性。

……………………………………

……………………………………

分析结果表明：在三个试验台站，雪深传感器与人工观测的积雪深度结果相一致，能够很好的反映出降雪和融雪期间雪面的变换，观测精度基本可以达到现有气象观测规范的要求。从测量结果的稳定性和准确性来看，激光雪深传感器的测量结果优于超声波雪深传感器。
\cuitkeywords{雪深，观测，传感器，对比分析}
\end{cuitabstractcn}

\begin{cuitabstracten}
This article is completed under the patronizing of China Meteorological Administration, public welfare meteorological profession specific (Project number: GYHY201006045). The paper through to the man-power and the ultrasonic wave snow depth sensor, the man-power and the laser snow depth sensor as well as the ultrasonic wave snow depth sensor and the laser snow depth sensor's tentative data contrastive analysis, has made system's theoretical analysis and the thorough performance analytical study to the ultrasonic wave snow depth sensor and the laser snow depth sensor's measuring technique.

This profession special to study the ultrasonic wave snow depth sensor and the laser snow depth sensor's survey characteristic, has designed the laboratory experiment and the outfield contrast observation experiment, the contrast ultrasonic wave snow depth sensor and the laser snow depth sensor's survey performance, simultaneously the data which surveys to the snow depth sensor has made system's analysis and the research, analyzes its data and the synoptic process (wind speed, temperature, snows with thaw process and so on) between relations as well as the snow depth and the snowfall relations. Summarizes the primary coverage which this article studies as follows to show:

(1) Designs the laboratory examination plan, through in the laboratory to sensor's static examination, obtains the data which the ultrasonic wave snow depth sensor and the laser snow depth sensor surveys, analyzes sensor's uniformity and the stability.

……………………

……………………
\cuitenglishkeywords{Snow Depth, Observation, Sensor, Contrastive Analysis}
\end{cuitabstracten}

\cuittableofcontents
\clearpage

\cuitmainmatter
\chapter{绪论}

\section{课题研究的重要意义}

本课题来源于2010年中国气象局公益性气象行业专项“冰雪自动化观测技术研究”（项目编号：GYHY201006045）。

我国是雨雪、冰冻灾害性天气多发的国家，具有影响范围广、程度深、发生频次高的特点。2008年初，我国南方地区连续遭受低温雨雪冰冻极端天气过程袭击，导致高速封路、机场关闭、高压线路破坏，对电力、通信、交通和人民生命财产带来极大危害。2009年末至2010年3月初，我国北方十多个省市地区遭受多次雪灾，给交通运输造成严重影响，对人们的生产生活特别是农林畜牧业生产造成巨大损失。

………………

………………

………………

………………

\section{雪深自动测量的研究现状}

\subsection{几种常用的雪深测量方法}

首先，介绍传统的人工测量雪深的方法，指出此方法的优点以及不足。然后介绍由此而产生的雪深自动测量的方法，主要都有哪些测量方法，简要介绍各种方法的特点\cite{ref-wetland,ref-water}。

………………

………………

光扫描法

投光器和受光器彼此保持一定距离，口朝下安装在距离地面的某一高度上。受光器的方向固定，投光器则常对雪面左右扫描，投光器的光进入受光器的视野时，根据投光器的角度求出积雪深度。

双杆法

将方向相对的投光器和受光器装在有数米间距的直立的两根杆子上，使其可沿杆同时上下滑动。当投光器、受光器位于雪面下方时，来自投光器的光受到阻挡，当它们一超出雪面，受光器就接收到投光器的光。根据这时的高度求出积雪深度。

………………

………………

表格格式：

\begin{table}[htbp]
  \centering
  \caption{“正确”的雪深瞬时值的判断条件\cite{ref-standard,ref-lantian}}
  \begin{tabular}{ccccc}
    \toprule
    气象变量 & 下限 & 上限 & 存疑的变化速率 & 错误的变化速率 \\
    \midrule
    雪深 & 0 mm & 1500 mm & 20--60 mm & $>60$ mm \\
    \bottomrule
  \end{tabular}
\end{table}

\begin{table}[htbp]
  \centering
  \caption{试验测试用具及性能指标}
  \begin{tabular}{cccc}
    \toprule
    序号 & 名称 & 型号 & 主要性能指标 \\
    \midrule
    1 & 传感器 & A型 & 分辨率1 mm \\
    2 & 采集器 & B型 & 采样间隔1 min \\
    \bottomrule
  \end{tabular}
\end{table}

图片格式（本图为例样，请根据论文情况据实画图）：

\begin{figure}[htbp]
  \centering
  \fbox{\rule{0pt}{6cm}\rule{8cm}{0pt}}
  \caption{试验装置图}
\end{figure}

公式格式：

\begin{equation}
  E = mc^2
\end{equation}

\setcounter{chapter}{4}
\chapter{结论与展望}

\section{结论}

根据雪深传感器测量特点，分别从仪器运行稳定性、观测数据的完整性、可靠性、准确性等角度对参试仪器性能进行综合评价，得到以下结论：

雪深自动化观测方法可行，雪深传感器组成结构和外观设计基本合理，在低温、大风、暴雪等恶劣环境下连续正常工作，未出现运行事故。参试雪深传感器与人工对比观测之间的准确性较好，可以连续、实时监测积雪深度的变化，基本可以满足冬季雪深自动化观测需求。

……………………

……………………

\section{展望}

为了进一步提高超声波传感器的测量精度，提出以下建议：

提高超声波传感器的耐寒强度，使其能在高寒环境下正常工作。

选择避风条件比较好的地点进行观测或采用防风措施。

……………………

……………………

根据试验数据分析结果，应该改进传感器采样算法和数据质量控制方法，提高传感器观测数据的可靠性和传感器运行稳定性。


\begin{thebibliography}{99}
\bibitem{ref-wetland} 雷光春.综合湿地管理：综合湿地管理国际研讨会论文集[C].北京:海洋出版社,2012.
\bibitem{ref-water} 马欢.人类活动影响下海河流域典型区水循环变化分析[D].北京:北京大学,2011.
\bibitem{ref-standard} 全国信息与文献标准化技术委员会.文献著录:第4部分非书资料:GB/T 3792.4-2009,[S].北京:中国标准出版社,2010.
\bibitem{ref-lantian} 袁训来,陈哲,肖书海,等.蓝田生物群：一个认识多细胞生物起源和早期演化的新窗口[J].科学通报,2012,55(34):3219.
\bibitem{ref-news} 丁文祥.数字革命与竞争国际化[N].中国青年报,2000-11-20(15).
\bibitem{ref-web} 萧钰.出版业信息化迈入快车道[EB/OL].(2001-12-19)[2002-04-15].\nolinkurl{http:www.creader.com/news.20011219/200112190019.html}.
\bibitem{ref-cimo} CIMO Survey on National Summaries of Methods and Instruments Related to Solid Precipitation Measurement at Automatic Weather Stations[J]. WMO/TD No.1544, 2010.
\end{thebibliography}

\cuitachievement
\noindent{\songti\bfseries\fontsize{12bp}{20bp}\selectfont 1 学术论文}\par
雷光春.综合湿地管理：综合湿地管理国际研讨会论文集[C].北京:海洋出版社,2016.（已录用）

XXX,XXX,XX,等.蓝田生物群：一个认识多细胞生物起源和早期演化的新窗口[J].科学通报,2012,55(34):3219.（已见刊）

\noindent{\songti\bfseries\fontsize{12bp}{20bp}\selectfont 2 科研项目}\par
[1]参与项目,XXX, XXX, XXX等. 我国西部地区冬季雪盖变化分析, 公益性行业（气象）科研专项（GYHYXXXXXXX）, 201X.XX.

\noindent{\songti\bfseries\fontsize{12bp}{20bp}\selectfont 3 专利与获奖}\par
[1]发明专利,XXX, XXX, XXX等. 一种用于雪盖遥感的数据处理方法（ZL201XXXXXXXX.X）, 201X.XX.

[2]第X届全国电子设计竞赛,XXX, XXX, XXX等. X等奖, 201X.XX.

\cuitacknowledgement
对导师(姓名、职称)的悉心指导表示感谢。

对授课教师(姓名、职称)授课与答疑解惑表示感谢。

对完成本文提供过帮助的人员（姓名、职称或学衔）表示感谢。

对同窗朋友（姓名）的帮助与支持表示感谢。

\end{document}
